% appendix.tex
% Supplementary material: empirical validations of assumptions used in
% the main-text derivations, plus material that is too dense for the
% main body but useful for replication.
%
% Language-audit pass, 2026-07-14 (see project_language_audit_2026-07-14
% memory; final section of the pass). Kept every algebra step (no
% collapsed align blocks) but cut narrating prose around them ("we
% expand this in two stages", "distributing", "pulling outside the
% sums"). Killed generic terminology: "artifact" -> "bias" (matches
% discussion.tex's "diversification-return bias"; the one surviving
% "sample-size artifact" is a standard statistical idiom, kept
% consistent with results.tex), "anchor"/"anchoring" -> literal
% ("set from", "matching"), "engine"/"engine-raw" -> "raw" (was
% implementation jargon for the uncorrected ensemble; also renamed the
% $W_{\mathrm{engine}}$ subscript to $W_{\mathrm{raw}}$), "operating
% point" -> "parameter values"/"fitted", "pipeline"/"framework" ->
% "construction"/"procedure", generic "mechanism" -> literal (e.g. "the
% persistence mechanism" -> "the duration model"), "downstream" cut,
% "smoke test" cut (software-testing jargon), "composer" -> "composition
% method" (last file needing this swap; done everywhere now), "sat/sits"
% -> "clustered"/"remained near"/"falls". Table 3 caption's "Clipping is
% a stability mechanism" rewritten to state the observed result
% directly, per the audit's own suggested rewrite.

\section{Empirical cross-term diagnostic}
\label{sec:supp-cov}

The variance correction of Eq.~\eqref{eq:scale} assumes
$\Cov(\tg_i^{\mathrm c},g_m)\approx0$, allowing the variance derivation to omit
the cross-term. Each generator draw is simulated independently of the
market path, so this covariance is zero in population, but on a
finite path its sample value is not. We measured the size of the
sample cross-term for three residual sources: the full-return
JumpHMM generator used by the hybrid and naive composition methods,
the residual-fit JumpHMM generator used by the JumpHMM-on-residuals
method, and the Gaussian residual draw of the Gaussian SIM. For each
of the $423$ non-market tickers, we sampled $R = 100$ training-length
growth-rate paths $\tg_i^{(r)}$ from each source and computed the
sample correlation of each centered draw with the observed training
growth-rate path $g_m$ of SPY, averaged over the $R$ draws:
\begin{equation}
  \mathrm{cov}^{\mathrm{norm}}_i
  = \frac{1}{R} \sum_{r=1}^{R}
    \frac{\Cov((\tg_i^{(r)})^{\mathrm c},\, g_m)}{\sigma_{\gen,i}^{(r)}\, \sigma_m},
\end{equation}
where $\sigma_{\gen,i}^{(r)}$ is the sample standard deviation of
draw $r$. Centering does not change covariance, so these values equal
those obtained from the raw draws.
We plotted the per-ticker values as histograms
(Fig.~\ref{fig:cov_diag}). The universe means were
$1.53\times10^{-4}$, $6.58\times10^{-5}$, and $2.26\times10^{-4}$ for
the full-return, residual-fit, and Gaussian sources, respectively,
and the per-ticker values of all three ranged from $-0.0061$ to
$0.0061$. For two independent series of $2{,}766$ observations, the
sample correlation has standard deviation near $1/\sqrt{2{,}766}
\approx 0.019$, the spread we observed across single draws, so an
average over $100$ draws has standard deviation near $0.002$. The
per-ticker values are therefore of the size expected from sampling
noise, and the omitted cross-term is negligible on average. On a
single path, it is not zero, which is why the composed variance
matches its target only approximately.

\begin{figure}[H]
  \centering
  \includegraphics[width=0.85\textwidth]{figs/supplement/FigS01-Cross-Covariance.pdf}
  \caption{\textbf{Empirical normalized cross-covariance
    $\Cov(\tg_i^{\mathrm c}, g_m) / (\sigma_{\gen,i}\, \sigma_m)$ across
    the $423$ non-market tickers.} Each per-ticker value is the sample
    correlation between a centered generator draw and the observed
    training SPY growth-rate path, averaged over $R = 100$ draws;
    the histograms for the three residual sources are stacked. Universe means were $1.53\times10^{-4}$,
    $6.58\times10^{-5}$, and $2.26\times10^{-4}$ for the full-return
    JumpHMM, residual-fit JumpHMM, and Gaussian generators.}
  \label{fig:cov_diag}
\end{figure}

% =============================================================================
% S2  Bias correction for daily-rebalanced backtests
% Four flat paragraphs:
%   (1) origin of the bias (Booth-Fama identity + iid residuals)
%   (2) empirical demonstration on a 20-asset allocator
%   (3) the location-shift correction (formula, prior, validation check)
%   (4) rejected alternatives
% =============================================================================

\section{Rebalancing under independent residuals}
\label{sec:supp-bias}

To determine what a rebalanced portfolio of composed paths
(Eq.~\eqref{eq:sim-composition}) inherits from the independent-residual
construction, this appendix derives the portfolio's diversification
return. On the variance-preserving branch, the variance correction in
Eq.~\eqref{eq:scale} sets each composed variance to the generator
variance. When the generator variance matches the empirical variance,
the per-asset variance terms of the diversification return derived
below are matched as well. What the construction does not match is
the cross-asset residual covariance, and the derivation isolates where
that omission enters.

The composition is built to reproduce per-asset marginals and
single-factor cross-sectional dependence. Each asset's residual path,
whether a rescaled full-return draw or a draw from a residual model,
keeps its own temporal structure, but the residual paths are generated
independently across assets at each date, so the construction omits
cross-asset residual covariance. Consider a daily-rebalanced portfolio
with constant target weights $w_i \ge 0$, $\sum_i w_i = 1$,
reset to those weights at the start of each trading day. Throughout
this derivation $g_i(t)$ is the annualized log-growth rate of
Eq.~\eqref{eq:growth_rate}, so the per-period log return on asset
$i$ is $g_i(t)\,\Delta t$ with $\Delta t = 1/252$~yr. The wealth
process is:
\begin{equation}
  W_t \;=\; W_{t-1} \sum_{i} w_i \exp\!\big(g_i(t)\,\Delta t\big),
  \label{eq:wealth-recursion}
\end{equation}
so the per-step log return on the portfolio is
$\log(W_t/W_{t-1}) = \log\!\sum_i w_i \exp(g_i(t)\,\Delta t)$. We
suppress the time argument, writing $g_i$ for $g_i(t)$. A
second-order Taylor expansion of the exponential about zero is
$\exp(g_i\,\Delta t) = 1 + g_i\,\Delta t + \tfrac{1}{2}(g_i\,\Delta t)^2
+ O((g\,\Delta t)^3)$. Substituting it into each summand of
$\sum_i w_i \exp(g_i\,\Delta t)$ gives:
\begin{align}
  \sum_{i} w_i \exp\!\big(g_i\,\Delta t\big)
   &= \sum_{i} w_i \!\left[\,1 + g_i\,\Delta t
        + \tfrac{1}{2}(g_i\,\Delta t)^2
        + O\!\big((g\,\Delta t)^3\big)\right]
       && \text{(Taylor)} \notag \\
   &= \sum_{i} w_i
        + \Delta t \sum_{i} w_i\, g_i
        + \tfrac{1}{2}(\Delta t)^2 \sum_{i} w_i\, g_i^2
        + O\!\big((g\,\Delta t)^3\big)
       && \text{(distribute)} \notag \\
   &= 1 + \Delta t \sum_{i} w_i\, g_i
        + \tfrac{1}{2}(\Delta t)^2 \sum_{i} w_i\, g_i^2
        + O\!\big((g\,\Delta t)^3\big).
       && \text{(}\textstyle\sum_i w_i = 1\text{)}
  \label{eq:exp-expand}
\end{align}
A second-order Taylor expansion of $\log(1+x)$ about $x=0$ gives
$\log(1+x) = x - \tfrac{1}{2}x^2 + O(x^3)$. Identifying $1 + x$ with
the right-hand side of Eq.~\eqref{eq:exp-expand}, the part beyond one
is
$x :=\Delta t \sum_i w_i\, g_i + \tfrac{1}{2}(\Delta t)^2 \sum_i w_i\, g_i^2$,
and squaring keeps only the leading $(\Delta t)^2$ term:
$x^2 = (\Delta t)^2 (\sum_i w_i\, g_i)^2 + O((\Delta t)^3)$ because
the cross term has an extra $\Delta t$ factor. Substitution gives:
\begin{align}
  \log\!\frac{W_t}{W_{t-1}}
   &= \log\!\Big(1 + x + O\!\big((g\,\Delta t)^3\big)\Big) \notag \\
   &= x - \tfrac{1}{2}\, x^2 + O\!\big(x^3\big) \notag \\
   &= \Delta t \sum_{i} w_i\, g_i
        + \tfrac{1}{2}(\Delta t)^2 \sum_{i} w_i\, g_i^2
        \notag \\
   &\quad - \tfrac{1}{2}(\Delta t)^2
        \!\Big(\sum_{i} w_i\, g_i\Big)^{\!2}
        + O\!\big((g\,\Delta t)^3\big) \notag \\
   &= \Delta t \sum_{i} w_i\, g_i
        + \tfrac{1}{2}(\Delta t)^2 \!\left[\sum_{i} w_i\, g_i^2
        \right. \notag \\
   &\hspace{8em}\left.{}- \Big(\sum_{i} w_i\, g_i\Big)^{\!2}\right]
        + O\!\big((g\,\Delta t)^3\big).
  \label{eq:rp-taylor}
\end{align}
Dividing through by $\Delta t$ extracts the portfolio's annualized
one-step log-growth rate $g_p(t) := \Delta t^{-1}\log(W_t/W_{t-1})$,
the portfolio counterpart of the per-asset $g_i(t)$ (still
suppressing the time argument):
\begin{equation}
  g_p \;=\; \sum_{i} w_i\, g_i
   \;+\; \tfrac{1}{2}\,\Delta t \!\left(\sum_{i} w_i\, g_i^2
        \;-\; \Big(\sum_{i} w_i\, g_i\Big)^{\!2}\right)
   \;+\; O\!\big((\Delta t)^2\big).
  \label{eq:gp-onestep}
\end{equation}

To obtain expected annualized portfolio growth, we restored the time
argument and took unconditional expectations of
Eq.~\eqref{eq:gp-onestep} under stationarity. Define the mean and
variance of each asset's annualized growth rate:
\begin{equation}
  \mu_i \;:=\; \E\!\big(g_i(t)\big), \qquad
  \sigma_i^2 \;:=\; \Var\!\big(g_i(t)\big),
\end{equation}
the covariance of the annualized rates
$\Sigma_{ij} := \Cov\!\big(g_i(t), g_j(t)\big)$, and the portfolio
variance
$\sigma_p^2 := \Var\!\big(\textstyle\sum_i w_i g_i(t)\big) =
w^\top \Sigma w$. The $O(\Delta t)$ correction in
Eq.~\eqref{eq:gp-onestep} contains two squared quantities,
$\sum_i w_i\, g_i(t)^2$ and $(\sum_i w_i\, g_i(t))^2$. Both expand
by the second-moment identity, which holds for any random variable
$X$ with finite second moment:
\begin{equation}
  \E\!\big(X^2\big) \;=\; \E\!\big(X\big)^2 \;+\; \Var(X).
  \label{eq:second-moment}
\end{equation}
Applying this identity to $X = g_i(t)$ and summing against $w_i$ gives:
\begin{align}
  \E\!\Big(\sum_{i} w_i\, g_i(t)^2\Big)
   &\;=\; \sum_{i} w_i\, \E\!\big(g_i(t)^2\big)
       && \text{(linearity)} \notag \\
   &\;=\; \sum_{i} w_i\!\left[\,\mu_i^2 + \sigma_i^2\,\right]
       && \text{(by Eq.~\eqref{eq:second-moment})} \notag \\
   &\;=\; \sum_{i} w_i\, \mu_i^2 \;+\; \sum_{i} w_i\, \sigma_i^2.
   \label{eq:exp-gi2}
\end{align}
Applying the identity to $X = \sum_i w_i g_i(t)$, whose mean is
$\sum_i w_i\mu_i$ and whose variance is $\sigma_p^2$, gives:
\begin{align}
  \E\!\Big(\Big(\textstyle\sum_{i} w_i g_i(t)\Big)^{\!2}\Big)
   &\;=\; \E\!\Big(\textstyle\sum_{i} w_i g_i(t)\Big)^{\!2}
       \;+\; \Var\!\Big(\textstyle\sum_{i} w_i g_i(t)\Big)
       && \text{(by Eq.~\eqref{eq:second-moment})} \notag \\
   &\;=\; \Big(\textstyle\sum_{i} w_i\, \mu_i\Big)^{\!2}
       \;+\; \sigma_p^2.
   \label{eq:exp-sum2}
\end{align}
We substituted Eqs.~\eqref{eq:exp-gi2}--\eqref{eq:exp-sum2} into the
expectation of Eq.~\eqref{eq:gp-onestep} and grouped the variance and
mean contributions as follows:
\begin{equation}
\begin{aligned}
  \E(g_p) \;={}& \sum_{i} w_i\, \mu_i
  + \underbrace{\tfrac{1}{2}\,\Delta t\!\left(\sum_{i} w_i\, \sigma_i^2
       - \sigma_p^2\right)}_{\displaystyle D} \\
  &+ \tfrac{1}{2}\,\Delta t\!\left[\sum_{i} w_i\, \mu_i^2
       - \Big(\sum_{i} w_i\, \mu_i\Big)^{\!2}\right]
  + O\!\big((\Delta t)^2\big),
\end{aligned}
  \label{eq:booth-fama}
\end{equation}
the Booth--Fama decomposition, to second order, of the growth of a
daily-rebalanced portfolio in annualized growth-rate units
\cite{boothFama1992, markowitz1976longRun}. The middle term, $D$, is
the diversification return: $\tfrac{1}{2}\Delta t$ times the
difference between the weighted-average per-asset variance
$\sum_i w_i \sigma_i^2$ and the portfolio variance $\sigma_p^2$. The
moments $\mu_i$ and $\sigma_i^2$ of the annualized rates
have units $\mathrm{yr}^{-1}$ and $\mathrm{yr}^{-2}$ respectively;
the products $\Delta t\,\sigma_i^2$ and $\Delta t\,\sigma_p^2$ in
Eq.~\eqref{eq:booth-fama} are the annualized variances
($\mathrm{yr}^{-1}$). The third term in Eq.~\eqref{eq:booth-fama} is
the cross-sectional dispersion of the per-asset means, of order
$\Delta t\,\mu^2$, while $D$ is of order $\Delta t\,\sigma^2$. For
SPY, $\mu^2$ lies more than two orders of magnitude below $\sigma^2$
(Table~\ref{tab:descriptive_stats}), and we drop the third term,
treating the same ordering as typical of daily equity growth rates.
The realized annualized compound growth rate is
$\bar g_p = (T\,\Delta t)^{-1}\log(W_T/W_0) = T^{-1}\sum_{t} g_p(t)$.
If the composed growth rates form a stationary ergodic process,
$\bar g_p$ converges to $\E(g_p)$ by the ergodic theorem. The long-run
limit is then fixed by the unconditional joint distribution of
$(g_i(t))_i$ at a single date and does not depend on the serial
structure of the residual draw. This limit argument treats the
residuals as a stationary process; the per-path centering and scaling
of the composition act on a finite horizon and are outside it.

To show that $D$ is non-negative, we rewrote the bracket in $D$ as a
sum over asset pairs. Expanding the portfolio variance, using
$1 - w_i = \sum_{j \ne i} w_j$ for weights summing to one, and
collecting each pair once gives:
\begin{align}
  \sum_i w_i \sigma_i^2 - \sigma_p^2
   &\;=\; \sum_i w_i \sigma_i^2 - \sum_i w_i^2 \sigma_i^2
     - 2 \sum_{i<j} w_i w_j \Sigma_{ij}
     && \text{(expand $\sigma_p^2$)} \notag \\
   &\;=\; \sum_i w_i (1 - w_i)\, \sigma_i^2
     - 2 \sum_{i<j} w_i w_j \Sigma_{ij}
     && \text{(collect)} \notag \\
   &\;=\; \sum_i \sum_{j \ne i} w_i w_j\, \sigma_i^2
     - 2 \sum_{i<j} w_i w_j \Sigma_{ij}
     && \text{($1 - w_i = \textstyle\sum_{j \ne i} w_j$)} \notag \\
   &\;=\; \sum_{i<j} w_i w_j
     \left(\sigma_i^2 + \sigma_j^2 - 2\Sigma_{ij}\right)
     && \text{(pair terms)} \notag \\
   &\;=\; \sum_{i<j} w_i w_j \,\Var(g_i - g_j),
     && \text{(variance of a difference)}
  \label{eq:pair-identity}
\end{align}
so that:
\begin{equation}
  D \;=\; \tfrac{1}{2}\,\Delta t \sum_{i<j} w_i w_j \,\Var(g_i - g_j)
  \;\ge\; 0.
  \label{eq:dr-pairwise}
\end{equation}
The diversification return is therefore $\tfrac{1}{2}\Delta t$ times
a sum of pairwise difference variances with pair weights $w_i w_j$.
It is zero only when every pair of assets with positive weight
comoves perfectly ($\Var(g_i - g_j) = 0$) or all weight is on one
asset. The rebalancing rule itself produces
this effect: weights drift between rebalances, the rebalancer sells
whatever has outperformed (its weight rose above target) and buys
whatever has underperformed (its weight fell below target), and the
trade systematically converts period-by-period cross-sectional
dispersion into compounded growth. Markowitz
\cite{markowitz1976longRun} described this as the variance bonus of
the long-run mean-variance allocation; Booth and Fama
\cite{boothFama1992} called it the diversification return.

To separate market and residual contributions to the diversification
return, we substituted the single-index decomposition
$g_i(t) = \alpha_i + \beta_i\, g_m(t) + \eps_i(t)$ into
Eq.~\eqref{eq:dr-pairwise}, with market variance
$\sigma_m^2 := \Var(g_m(t))$, per-asset residual variance
$\sigma_{\eps,i}^2 := \Var(\eps_i(t))$, and $\Cov(g_m, \eps_i) = 0$
(checked empirically in Online Appendix~\ref{sec:supp-cov}). The
pairwise difference variance becomes:
\begin{equation}
  \Var(g_i - g_j) \;=\; (\beta_i - \beta_j)^2 \sigma_m^2
  \;+\; \Var(\eps_i - \eps_j).
\end{equation}
Because $g_m$ and the residuals are uncorrelated, this relation splits
the diversification return into market-beta-dispersion and
idiosyncratic components:
\begin{equation}
  D \;=\;
  \underbrace{\tfrac{1}{2}\,\Delta t\,\sigma_m^2 \sum_{i<j} w_i w_j (\beta_i - \beta_j)^2}_{D_{\mathrm{mkt}}}
  \;+\;
  \underbrace{\tfrac{1}{2}\,\Delta t \sum_{i<j} w_i w_j \,\Var(\eps_i - \eps_j)}_{D_\eps},
  \label{eq:dr-sim}
\end{equation}
where the market-beta-dispersion component $D_\mathrm{mkt}$ contains
the weighted cross-sectional variance of $\beta$. Applying the pair
identity of Eq.~\eqref{eq:pair-identity} with $\beta_i$ in place of
$g_i$ gives:
\begin{align}
  \sum_{i<j} w_i w_j (\beta_i - \beta_j)^2
   &\;=\; \sum_i w_i (1 - w_i)\, \beta_i^2
     - 2 \sum_{i<j} w_i w_j \beta_i \beta_j
     && \text{(pair identity)} \notag \\
   &\;=\; \sum_i w_i \beta_i^2
     - \Big(\sum_i w_i^2 \beta_i^2
     + 2 \sum_{i<j} w_i w_j \beta_i \beta_j\Big)
     && \text{(regroup)} \notag \\
   &\;=\; \sum_i w_i \beta_i^2 - \Big(\sum_i w_i \beta_i\Big)^{\!2},
     && \text{(square of a sum)}
\end{align}
so $D_\mathrm{mkt} = \tfrac{1}{2}\,\Delta t\,\sigma_m^2 [\sum_i w_i
\beta_i^2 - (\sum_i w_i \beta_i)^2]$, the weighted variance of $\beta$
scaled by the annualized market variance $\Delta t\,\sigma_m^2$.
Across the $423$ non-market tickers, the calibrated $\beta$
distribution spans $[0.03, 1.79]$ (Fig.~\ref{fig:branch_map}). A
weighted variance is at most one quarter of the squared range, so
$D_{\mathrm{mkt}} \le \tfrac{1}{8}\,\Delta t\,\sigma_m^2\,(1.79 -
0.03)^2$. The SPY closing-price series used for the multi-asset
universe has annualized variance $\Delta t\,\sigma_m^2 =
0.030~\mathrm{yr}^{-1}$ (Fig.~\ref{fig:preservation}), at which
$D_{\mathrm{mkt}}$ is at most $0.012~\mathrm{yr}^{-1}$, or $1.2$
percentage points per year. For any long-only weights on this
universe, the market component therefore contributes little, and a
larger diversification return must come from the idiosyncratic
component $D_\eps$. In the synthetic construction, residual paths are
generated independently across assets, so
$\Var(\eps_i - \eps_j) = \sigma_{\eps,i}^2 + \sigma_{\eps,j}^2$ and,
by the same pair identity,
$D_\eps = \tfrac{1}{2}\,\Delta t \sum_i w_i (1 - w_i)\,
\sigma_{\eps,i}^2$, a weighted sum of the annualized idiosyncratic
variances $\Delta t\,\sigma_{\eps,i}^2$.

Because the second-order value of $D$ depends only on same-date
unconditional moments, serial dependence in the residual draw does not
change it. Equation~\eqref{eq:scale} chooses $s_i$ so that each
composed variance $\beta_i^2\sigma_m^2+\sigma_{\eps,i}^2$ equals the
generator variance $\sigma_{\gen,i}^2$, and any residual generator
with the same same-date variances gives the same $D$. As an
illustration, consider a stationary first-order autoregressive
[AR(1)] residual $\eps_i(t) = \phi\, \eps_i(t-1) + \eta_i(t)$ with
independent and identically distributed innovations $\eta_i$ whose
variance is chosen so that
$\Var(\eps_i) = \sigma_{\eta,i}^2 / (1 - \phi^2) = \sigma_{\eps,i}^2$.
The parameter $\phi$ leaves the variance at a single date unchanged,
so it does not change $\E(g_p)$ in Eqs.~\eqref{eq:booth-fama}
and~\eqref{eq:dr-sim}. It does change sampling noise: for an
uncentered stationary AR(1) residual with variance $\sigma_\eps^2$,
the variance of the sample mean is approximately
$T^{-1}\sigma_\eps^2(1+\phi)/(1-\phi)$. Serial dependence therefore
changes the per-path noise around the stationary mean, not the mean
itself, and is not a lever on the diversification return at the level
of the second-order formula.

The difference between synthetic and real residuals that does change
$D$ is contemporaneous. The derivation assumes residuals are
uncorrelated across assets at each date, as they are in the
independent per-asset generator. Real residuals share sector and style
factors, which reduce $\Var(\eps_i-\eps_j)$ below the
independent-generator value $\sigma_{\eps,i}^2+\sigma_{\eps,j}^2$.
Matching each marginal variance does not restore this covariance, so
the synthetic $D_\eps$ exceeds its real-market counterpart when the
omitted covariance is positive. We did not measure the size of this
excess, which would require a joint-residual experiment with the
marginals and the allocation rule held fixed. A comparison with
realized portfolio returns would also contain a second effect:
realized diversification returns on real equity portfolios fall short
of the second-order formula at all rebalancing
horizons~\cite{boucheyNemtchinovWong2015}. A measured excess of
synthetic over realized diversification return would therefore bound
the part caused by independent residuals from above rather than
measure it.


% =============================================================================
% S3 Validation metric definitions (KS, AD, W1, Hellinger)
% =============================================================================

\section{Validation metric definitions}
\label{sec:supp-validation}

To distinguish marginal agreement from temporal agreement, we used
separate distributional, tail, and autocorrelation summaries
(Tables~\ref{tab:model_comparison} and~\ref{tab:jump_ablation}).
This section defines those summaries; the factor-recovery errors and
variance ratios of the composition scorecard are defined in the main
text with that experiment. For one observed series $x=(x_1,\ldots,x_n)$ and $R$ synthetic
paths $y^{(r)}$ of the same length, each distributional comparison
used one complete synthetic path and the fixed observed history.
The two-sample Kolmogorov--Smirnov (KS) statistic is the largest
absolute difference between their empirical cumulative distribution
functions, $D_{\mathrm{KS}}=\sup_z|F_x(z)-F_{y^{(r)}}(z)|$
\cite{kolmogorov1933,smirnov1948}. The Anderson--Darling (AD)
comparison gives greater weight to differences near the tails
\cite{andersonDarling1952}. We used the approximate two-sample KS
test and the modified, tie-aware two-sample AD calculation in
\texttt{HypothesisTests.jl}, with asymptotic $p$-values. A path
passed when $p>0.05$, and the pass percentage was $100$ times the
fraction of passing paths. These percentages are not calibrated
significance levels, for three reasons: returns are serially
dependent, whereas the tests assume independent observations; every
path is compared with the same observed series, so the comparisons
are not independent of one another; and in sample, the generator was
fitted to that series. A high pass rate therefore describes agreement
with the supplied history and is not evidence that the model and the
data share a distribution.

To quantify marginal discrepancies without a pass threshold, we
computed Wasserstein-1 and histogram Hellinger scores for each path
(Table~\ref{tab:model_comparison}). Let $x_{(j)}$ and $y^{(r)}_{(j)}$
denote the $j$th order statistics in increasing order. Because each
synthetic path has the same length as the observed series, the
Wasserstein-1 distance is:
\begin{equation}
 W_1(x,y^{(r)})=\frac{1}{n}\sum_{j=1}^{n}
       |x_{(j)}-y^{(r)}_{(j)}|.
 \label{eq:wasserstein}
\end{equation}
It has the growth-rate units $\mathrm{yr}^{-1}$. For each Hellinger
comparison, we formed $50$ equally spaced bins between the minimum
and maximum of the pooled pair of samples. Let $p_b$ and $q_b$ be
the observed and simulated bin counts divided by the sample length
$n$. The histogram score is:
\begin{equation}
 H(x,y^{(r)})=\left[\frac{1}{2}\sum_{b=1}^{50}
          (\sqrt{p_b}-\sqrt{q_b})^2\right]^{1/2}.
 \label{eq:hellinger}
\end{equation}
Every bin, including the last, contained its lower edge and excluded
its upper edge, so observations equal to the pooled maximum fell
outside all bins and that sample's bin weights sum to slightly less
than one; we state this convention because the reported scores were
computed with it. The score is a histogram approximation to the
Hellinger distance between the two densities. Its bin edges depend on
each pair, and its bins hold fewer observations in shorter samples, so
it is best compared across generators within one window; holdout
scores on $249$ observations run larger than training scores partly
for that reason. Smaller Wasserstein and Hellinger scores indicate
closer marginal agreement.

To describe tail shape, we computed excess kurtosis for each
synthetic path and compared its mean over paths with the observed value
(Table~\ref{tab:model_comparison}); the composition tables report
medians instead. With
$m_p(x)=n^{-1}\sum_{t=1}^{n}(x_t-\bar x)^p$, the estimator was
$\kappa(x)=m_4(x)/m_2(x)^2-3$, without a finite-sample bias
correction. The composition scorecard also used the Hill index of
absolute returns (Table~\ref{tab:aggregate}). Let
$X_{(1)}\ge\cdots\ge X_{(n)}$ be decreasingly ordered absolute
returns and $k=\max\{2,\lfloor0.05n\rfloor\}$. The implemented
upper-tail estimator is:
\begin{equation}
 \hat\xi=\frac{1}{k-1}\sum_{j=1}^{k-1}
             \log\!\left(\frac{X_{(j)}}{X_{(k)}}\right).
 \label{eq:hill}
\end{equation}
The estimate $\hat\xi$ is the reciprocal of the conventional tail
index, so larger values indicate a heavier absolute-return tail; the
threshold $X_{(k)}$ must be positive. These shape summaries do not
measure how often returns exceed a Value-at-Risk threshold.

To assess whether the simulated distributions bracketed the observed
one, we measured agreement at $99$ quantile levels
(Table~\ref{tab:model_comparison}). For $u_j=j/100$, $j=1,\ldots,99$,
let $Q_x(u_j)$ be the observed sample quantile and
$Q_r(u_j)$ the corresponding quantile of path $r$. Let
$L_j$ and $U_j$ be the $5$th and $95$th percentiles of
$\{Q_r(u_j):r=1,\ldots,R\}$. All quantiles used linear interpolation
between order statistics, the common software default (type~7). The reported
quantile-envelope coverage is:
\begin{equation}
 C=\frac{100}{99}\sum_{j=1}^{99}
       \mathbf{1}\{L_j\le Q_x(u_j)\le U_j\}.
 \label{eq:quantile-envelope-coverage}
\end{equation}
Larger values mean that more observed quantiles fall inside the
pointwise simulation envelopes. This statistic assesses marginal
quantile agreement; it does not estimate daily Value-at-Risk
exceedance coverage or simultaneous confidence coverage.

To quantify persistence in return magnitude, we used the sample
autocorrelation function (ACF) of absolute returns. Each series was
centered by its own full-series absolute-return mean,
$\bar a_x=n^{-1}\sum_t|x_t|$. At lag $\ell$, the calculation was:
\begin{equation}
 a_\ell(x)=
 \frac{\sum_{t=1}^{n-\ell}(|x_t|-\bar a_x)(|x_{t+\ell}|-\bar a_x)}
      {\sum_{t=1}^{n}(|x_t|-\bar a_x)^2}.
 \label{eq:sample-absolute-acf}
\end{equation}
The denominator used all $n$ observations, with no lag-dependent
rescaling. Both error scores below start at lag one.
The single-asset comparison measured the error of the ensemble-average
ACF (Table~\ref{tab:model_comparison}). For a lag limit $L$, that
mean absolute error is:
\begin{equation}
 E^{\rm ensemble}_L=\frac{1}{L}\sum_{\ell=1}^{L}
 \left|a_\ell(x)-\frac{1}{R}\sum_{r=1}^{R}a_\ell(y^{(r)})\right|.
 \label{eq:acf-ensemble-error}
\end{equation}
We used $R=1{,}000$ and $L=\min(252,n-1)$: lags $1$--$252$
for the $2{,}766$-observation training series and $1$--$248$ for the
$249$-observation holdout. This evaluation lag limit is separate from
the $25$-lag window of the tuning objective in
Eq.~\eqref{eq:grid_search}. The state-resolution and emission
comparisons (Tables~\ref{tab:sensitivity_n}
and~\ref{tab:emission_hsmm}) used the same ensemble convention.

The multi-asset jump comparison instead measured error within each
path before averaging (Table~\ref{tab:jump_ablation}). For one
asset, its mean path error is:
\begin{equation}
 E^{\rm path}_L=\frac{1}{R}\sum_{r=1}^{R}
       \left[\frac{1}{L}\sum_{\ell=1}^{L}
       |a_\ell(x)-a_\ell(y^{(r)})|\right].
 \label{eq:acf-path-error}
\end{equation}
We used $L=25$ and $60$ in both windows and averaged these scores
over assets. The fallback sensitivity check also reported this path
score with $L=25$ (Table~\ref{tab:fallback_sensitivity}). At a common
lag limit and on the same paths, the triangle inequality gives
$E^{\rm ensemble}_L\le E^{\rm path}_L$: averaging the curves can
cancel errors of opposite sign. The ensemble score asks whether the
model's average autocorrelation curve matches the observed curve; the
path score asks how far a single path's curve typically falls from
it, sampling noise included. Their magnitudes
cannot be compared directly across the single-asset and composition
experiments, which also differ in lag limits, histories, and assets.

To describe simulation uncertainty in the single-asset comparison,
we used a standard-error method suited to each statistic
(Table~\ref{tab:model_comparison}). For a pass fraction $\hat p$,
the reported binomial standard error was
$100\sqrt{\hat p(1-\hat p)/R}$ percentage points. This estimate is
zero when every path passes or fails; it does not establish that the
underlying pass probability is exactly zero or one. For kurtosis,
Wasserstein distance, and Hellinger score, the standard error was
the sample standard deviation of the $R$ path scores divided by
$\sqrt R$. For ensemble ACF error and quantile-envelope coverage,
we used $500$ bootstrap resamples of the $R$ path indices, each of
size $R$ with replacement. Each resample retained entire paths and
recomputed the statistic across all lags or quantile levels. The
standard error was the sample standard deviation of those $500$
recomputed values, without a further division by $\sqrt{500}$.
All these errors condition on one observed history and a fixed fitted
model; they omit uncertainty in fitting, model selection, and market
regimes. The multi-asset jump contrasts used the replication as the
unit instead: within each of the $1{,}000$ replications we averaged a
metric across assets, differenced two jump settings, and took the
standard deviation of the $1{,}000$ differences divided by
$\sqrt{1{,}000}$ as the standard error, so that the dependence created
by the shared market path stays inside each unit
(Table~\ref{tab:jump_ablation_uncertainty}; Online
Appendix~\ref{sec:supp-jump-composition}).

% --- Figure: Statistical validation of HMM-WJ ---
\begin{figure}[H]
  \centering
  \includegraphics[width=\textwidth]{figs/supplement/FigS04-Statistical-Validation.pdf}
  \caption{\textbf{Statistical validation of the hidden Markov model with
    jumps (HMM-WJ)} ($N = 100$,
    $\alpha = 0.05$). Panels (a) and (d) condition on the
    $\sim\!24\%$ of in-sample paths that contain at least one jump
    episode. (a)~In-sample
    KS $p$-values for these $239$ paths (pass rate $93.7\%$). (b)~Out-of-sample
    KS $p$-values for all $1{,}000$ paths (pass rate $95.4\%$), since few
    holdout-length paths contain an episode. (c)~Out-of-sample
    density fan chart; the observed density falls within the simulation
    envelope. (d)~Observed and mean simulated in-sample autocorrelation
    function (ACF) of $|G_t|$
    for the paths with an episode.}
  \label{fig:statistical_validation}
\end{figure}


